% =====================================================================
%  tracks4u - Pitch de Ingeniería de Datos (Beamer + TikZ)
%  Overleaf: Menu > Compiler > pdfLaTeX. Compilar 2 veces.
% =====================================================================
\documentclass[aspectratio=169,11pt]{beamer}

\usepackage[utf8]{inputenc}
\usepackage[T1]{fontenc}
\usepackage[spanish,es-noshorthands]{babel}
\usepackage{lmodern}
\usepackage{tikz}
\usepackage{booktabs}
\usepackage{array}
\usepackage{tabularx}
\usepackage{adjustbox}
\usetikzlibrary{babel,shapes.geometric,shapes.multipart,arrows.meta,positioning,calc,fit,backgrounds}

% ---------- Paleta ----------
\definecolor{ink}{HTML}{0F172A}
\definecolor{muted}{HTML}{64748B}
\definecolor{line}{HTML}{CBD5E1}
\definecolor{soft}{HTML}{F8FAFC}
\definecolor{accent}{HTML}{FF5500}
\definecolor{batch}{HTML}{2563EB}
\definecolor{sysgray}{HTML}{475569}   % metadatos de sistema
\definecolor{wmgreen}{HTML}{16A34A}   % watermark / outbox
\definecolor{rawc}{HTML}{64748B}

% ---------- Tema minimalista ----------
\usetheme{default}
\setbeamertemplate{navigation symbols}{}
\setbeamersize{text margin left=9mm,text margin right=9mm}
\setbeamercolor{normal text}{fg=ink}
\setbeamercolor{structure}{fg=accent}
\setbeamertemplate{itemize item}{\textcolor{accent}{\rule[0.25ex]{0.7ex}{0.7ex}}}
\setbeamertemplate{frametitle}{%
  \vskip0.45em
  {\color{accent}\scriptsize\bfseries\insertsectionhead}\par
  \vskip0.05em
  {\color{ink}\large\bfseries\insertframetitle}\par
  \vskip0.2em
  {\color{line}\rule{\textwidth}{0.6pt}}%
}
\setbeamertemplate{footline}{%
  \hbox to \paperwidth{\hspace{9mm}\color{muted}\tiny tracks4u $\cdot$ Ingeniería de Datos (Django + RDS + S3 + Glue + Kafka + ML)
  \hfill\insertframenumber\,/\,\inserttotalframenumber\hspace{9mm}}%
  \vskip2.5mm}

% ---------- Helpers ----------
\newcommand{\kick}[1]{{\color{accent}\footnotesize\bfseries\MakeUppercase{#1}}\par}
\newcommand{\pillS}{\tikz[baseline=(p.base)]\node[fill=accent!12,text=accent,rounded corners=3pt,inner sep=2pt,font=\scriptsize\bfseries](p){STREAMING};}
\newcommand{\pillB}{\tikz[baseline=(p.base)]\node[fill=batch!12,text=batch,rounded corners=3pt,inner sep=2pt,font=\scriptsize\bfseries](p){BATCH / CDC};}
\newcommand{\pillM}{\tikz[baseline=(p.base)]\node[fill=wmgreen!14,text=wmgreen,rounded corners=3pt,inner sep=2pt,font=\scriptsize\bfseries](p){ML / ANALYTICS};}
\newcommand{\PK}{\textcolor{accent}{PK}~}
\newcommand{\FK}{\textcolor{batch}{FK}~}

\tikzset{
  card/.style={draw=line,rounded corners=5pt,fill=white,inner sep=6pt},
  cardsoft/.style={draw=line,rounded corners=5pt,fill=soft,inner sep=6pt},
  flow/.style={-{Stealth[length=2mm]},line width=0.8pt,draw=ink},
  flowS/.style={-{Stealth[length=2mm]},line width=1pt,draw=accent},
  flowB/.style={-{Stealth[length=2mm]},line width=1pt,draw=batch},
  % --- diagrama de clases ---
  umlc/.style={rectangle split,rectangle split parts=3,draw=ink,line width=0.6pt,
    text width=29mm,align=left,inner sep=2.0pt,anchor=north,text=ink,
    font=\fontsize{5.2}{6.0}\selectfont\ttfamily},
  umlw/.style={rectangle split,rectangle split parts=2,draw=wmgreen,line width=0.6pt,
    text width=29mm,align=left,inner sep=2.0pt,anchor=north,text=ink,
    font=\fontsize{5.2}{6.0}\selectfont\ttfamily,
    rectangle split part fill={wmgreen,white}},
  umlo/.style={rectangle split,rectangle split parts=2,draw=accent,line width=0.6pt,
    text width=29mm,align=left,inner sep=2.0pt,anchor=north,text=ink,
    font=\fontsize{5.2}{6.0}\selectfont\ttfamily,
    rectangle split part fill={accent,white}},
  rel/.style={line width=0.6pt,draw=ink},
  mult/.style={font=\fontsize{5}{6}\selectfont\bfseries,text=ink,inner sep=1pt},
  verb/.style={font=\fontsize{5}{6}\selectfont\itshape,text=muted,inner sep=1.5pt},
}
% \umlclass{id}{(x,y)}{NOMBRE}{color header}{atributos}
\newcommand{\umlclass}[5]{%
  \node[umlc,rectangle split part fill={#4,white,sysgray}] (#1) at #2
  {\textcolor{white}{\rmfamily\bfseries #3}\nodepart{two}#5\nodepart{three}\textcolor{white}{trig: updated\_at}};}

\begin{document}

% ===================== 1 · PORTADA =====================
\begin{frame}[plain]
  \vfill
  \kick{Sistemas Intensivos en Datos $\cdot$ Arquitectura Implementada}
  \vspace{0.3em}
  {\color{ink}\fontsize{40}{44}\selectfont\bfseries tracks4u}\par
  \vspace{0.3em}
  {\color{muted}\Large De portal transaccional en Django a plataforma intensiva en datos en AWS}\par
  \vspace{0.6em}{\color{accent}\rule{3cm}{2.5pt}}\par
  \vspace{0.6em}
  {\color{ink}\normalsize Estado en Amazon RDS PostgreSQL $\cdot$ Multimedia y Data Lake en Amazon S3\\
  Ingesta Batch/CDC con AWS Glue $\cdot$ Streaming con Outbox y Apache Kafka $\cdot$ Recomendador Item-Item}
  \vfill
\end{frame}

% ===================== 2 · QUÉ ES TRACKS4U =====================
\section{¿Qué es tracks4u?}
\begin{frame}{Un portal de música instrumentado de punta a punta}
  \small
  \textcolor{muted}{tracks4u es una plataforma web en Django desplegada sobre AWS EC2, RDS PostgreSQL y S3. Los oyentes reproducen canciones servidas mediante URLs prefirmadas desde S3, arman playlists y compran canciones directamente.}

  \vspace{0.9em}
  \footnotesize
  \begin{columns}[T,onlytextwidth]
    \begin{column}{0.31\textwidth}
      \textbf{\color{accent}Escuchar y descubrir}
      \begin{itemize}\setlength{\itemsep}{2pt}
        \item Reproductor HTML5 instrumentado (\texttt{player.js} + \texttt{telemetry.js}).
        \item Catálogo de 36 canciones (\texttt{.opus}) y 33 portadas en S3 con duración extraída vía \texttt{mutagen}.
        \item Búsqueda por título, subgénero, artista y catálogo de 20 géneros.
      \end{itemize}
    \end{column}
    \begin{column}{0.31\textwidth}
      \textbf{\color{accent}Organizar y compartir}
      \begin{itemize}\setlength{\itemsep}{2pt}
        \item Playlists públicas/privadas con \texttt{tags} (\texttt{TEXT[]}) y reordenamiento transaccional.
        \item Eventos en vivo de \texttt{like}, \texttt{follow} y \texttt{share}.
        \item Emisión transaccional de \texttt{playlist\_add} desde el servidor.
      \end{itemize}
    \end{column}
    \begin{column}{0.31\textwidth}
      \textbf{\color{accent}Comprar y publicar}
      \begin{itemize}\setlength{\itemsep}{2pt}
        \item Compra directa con órdenes (\texttt{order} + \texttt{order\_item}) y precio congelado (\texttt{price\_paid}).
        \item Onboarding de artistas y subida validada a S3.
        \item Dashboard de analítica e ingresos desglosados por género.
      \end{itemize}
    \end{column}
  \end{columns}

  \vfill
  \begin{center}\itshape\textcolor{muted}{El estado vive en RDS PostgreSQL, los binarios en S3 y cada interacción viaja al Data Lake.}\end{center}
\end{frame}

% ===================== 3 · FUNCIONALIDADES IMPLEMENTADAS =====================
\section{Funcionalidades implementadas}
\begin{frame}{Qué construimos en la plataforma de datos}
  \footnotesize
  \begin{columns}[T,onlytextwidth]
    \begin{column}{0.48\textwidth}
      \textbf{\color{accent}Captura y Almacenamiento}
      \begin{itemize}\setlength{\itemsep}{2pt}
        \item \textbf{RDS PostgreSQL (\texttt{schema.sql}):} 12 tablas con triggers \texttt{trg\_<tabla>\_set\_updated\_at} y privilegios segregados (\texttt{grants.sql}).
        \item \textbf{Patrón Outbox (\texttt{telemetry\_event}):} ingesta por lotes en \texttt{POST /api/events/} y drenaje seguro con ventana de asentamiento de 5\,s.
        \item \textbf{Streaming Kafka:} \texttt{producer.py} drena el outbox hacia el topic \texttt{tracks4u.events} y confirma (\texttt{/ack/}) tras \texttt{flush()}.
      \end{itemize}
      \vspace{0.5em}
      \textbf{\color{accent}Ingesta al Data Lake (S3)}
      \begin{itemize}\setlength{\itemsep}{2pt}
        \item \textbf{Consumer Kafka $\rightarrow$ RAW:} deduplica por \texttt{event\_id}, escribe CSV y solo hace \texttt{commit()} del offset tras subir a S3.
        \item \textbf{AWS Glue Batch/CDC:} lectura con \texttt{etl\_watermark} y escritura particionada en S3.
      \end{itemize}
    \end{column}
    \begin{column}{0.48\textwidth}
      \textbf{\color{accent}Perfilado Conductual (\texttt{ANALYTICS})}
      \begin{itemize}\setlength{\itemsep}{2pt}
        \item Construcción automática de \texttt{entity=user\_profile} en Parquet con conteos por tipo de evento y cálculo de \texttt{skip\_rate}.
        \item Limpieza resiliente: extracción de \texttt{track\_id} desde \texttt{payload}, normalización de IDs y filtrado de anónimos.
      \end{itemize}
      \vspace{0.5em}
      \textbf{\color{accent}Recomendador ML Item-Item (\texttt{recommend.py})}
      \begin{itemize}\setlength{\itemsep}{2pt}
        \item Filtrado colaborativo Item-Item por similitud coseno sobre matriz ponderada usuario $\times$ canción.
        \item Validación \textit{leave-one-out} (\textbf{Hit-Rate@10 = 0.205} vs \textbf{0.030} popular y \textbf{0.033} azar).
        \item Despliegue aislado en Docker Compose (\texttt{network\_mode: host} + IAM Role \texttt{LabRole}).
      \end{itemize}
    \end{column}
  \end{columns}
\end{frame}

% ===================== 4 · DATOS BATCH =====================
\section{Datos en batch}
\begin{frame}{\pillB\ \ Esquema relacional en RDS PostgreSQL extraído por AWS Glue}
  \scriptsize
  \renewcommand{\arraystretch}{1.24}
  \begin{tabularx}{\textwidth}{@{}>{\bfseries\ttfamily}p{31mm} X >{\raggedright\arraybackslash}p{32mm}@{}}
    \toprule
    \textcolor{muted}{\rmfamily\mdseries Tabla en RDS} & \textcolor{muted}{Columnas y reglas implementadas} & \textcolor{muted}{Estrategia Glue}\\
    \midrule
    track\newline genre / track\_genre & title, subgenre, mood\_tags (\texttt{TEXT[]}), language, bpm, track\_type, price, duration\_sec, keys S3 + catálogo de 20 géneros & Snapshot diario / Incremental (\texttt{updated\_at})\\
    "user" & id, username, country, birth\_date, gender, ethnicity, created\_at, updated\_at (\textbf{sin} exponer email ni password) & Snapshot diario (columnas seguras)\\
    artist\newline artist\_genre & stage\_name, bio, country, user\_id (opcional para artistas de catálogo) + relación M2M con \texttt{genre} & Snapshot diario\\
    playlist\newline playlist\_track & name, description, tags (\texttt{TEXT[]}), cover\_image, is\_public; position con \texttt{UNIQUE DEFERRABLE} & Snapshot diario\\
    "order"\newline order\_item & total\_amount, currency (USD), payment\_method, status (\texttt{COMPLETED}); \texttt{price\_paid} congelado & \textbf{Incremental CDC} por \texttt{updated\_at}\\
    \bottomrule
  \end{tabularx}

  \vspace{0.6em}
  \textcolor{muted}{\textbf{Control CDC en RDS:} triggers \texttt{trg\_<tabla>\_set\_updated\_at} garantizan el avance de \texttt{updated\_at} incluso en cargas SQL; \texttt{etl\_watermark} registra \texttt{last\_value}, \texttt{rows\_loaded} y \texttt{status} por entidad.}
\end{frame}

% ===================== 5 · DATOS STREAMING =====================
\section{Datos en streaming}
\begin{frame}{\pillS\ \ Outbox \texttt{telemetry\_event} $\rightarrow$ Kafka (\texttt{tracks4u.events}) $\rightarrow$ S3 RAW}
  \scriptsize
  \renewcommand{\arraystretch}{1.18}
  \begin{tabularx}{\textwidth}{@{}>{\ttfamily}p{29mm} X >{\raggedright\arraybackslash}p{16mm} >{\raggedright\arraybackslash}p{36mm}@{}}
    \toprule
    \textcolor{muted}{\rmfamily Evento} & \textcolor{muted}{Payload validado (\texttt{telemetry/schemas.py})} & \textcolor{muted}{Peso ML} & \textcolor{muted}{Señal de negocio}\\
    \midrule
    play\_start        & track\_id, playlist\_id?                        & \textbf{+1} & Inicio de escucha\\
    play\_progress     & track\_id, position\_sec, percent\_listened     & —           & Retención intra-canción\\
    play\_complete     & track\_id, seconds\_listened                    & \textbf{+3} & Escucha completa\\
    skip               & track\_id, position\_sec, percent\_listened     & \textbf{-2} & Rechazo / cálculo skip\_rate\\
    pause / resume / seek & track\_id, position\_sec / from\_sec, to\_sec & —          & Interacción en reproductor\\
    like / unlike      & track\_id                                       & \textbf{+4} & Preferencia explícita\\
    follow / unfollow / share & artist\_id / track\_id, channel?         & —           & Viralidad y afinidad\\
    playlist\_add      & playlist\_id, track\_id, position (servidor)    & \textbf{+4} & Curaduría y co-ocurrencia\\
    search             & search\_id (UUID), query, results\_count, clicked\_track\_id? & — & Intención y CTR de búsqueda\\
    purchase\_intent   & track\_id, stage (\texttt{checkout\_start|completed|failed}), order\_id?, price? & \textbf{+5} & Conversión monetaria\\
    page\_view / dwell\_time & view, referrer? / seconds, max\_scroll\_pct? & —        & Engagement de sesión\\
    \bottomrule
  \end{tabularx}

  \vspace{0.5em}
  \textcolor{muted}{Columnas del CSV en \texttt{s3://.../raw/entity=play\_events/ingest\_date=AAAA-MM-DD/}: \texttt{event\_id, ts, session\_id, user\_id, event\_type, track\_id, device, geo, entrypoint, payload}.}
\end{frame}

% ===================== 6 · DIAGRAMA DE CLASES =====================
\section{Modelo de datos}
\begin{frame}{Esquema relacional implementado en RDS (\texttt{database/schema.sql})}
  \centering
  \begin{adjustbox}{max width=\textwidth,max totalheight=0.75\textheight,center}
  \begin{tikzpicture}
    % ---- fila superior ----
    \umlclass{artist}{(0,0)}{ARTIST}{ink}{\PK id (BIGINT)\\ \FK user\_id (null)\\ stage\_name, bio\\ country\\ created\_at, updated\_at}
    \umlclass{track}{(3.8,0)}{TRACK}{ink}{\PK id (BIGINT)\\ \FK artist\_id\\ title, subgenre\\ language, mood\_tags[]\\ bpm, track\_type, price\\ audio\_file, cover\_image\\ duration\_sec, updated\_at}
    \umlclass{genre}{(7.6,0)}{GENRE + PUENTES}{ink}{\PK id, name (20 base)\\ --- track\_genre ---\\ \FK track\_id, \FK genre\_id\\ --- artist\_genre ---\\ \FK artist\_id, \FK genre\_id}
    \umlclass{plt}{(11.4,0)}{PLAYLIST\_TRACK}{accent}{\PK id (BIGINT)\\ \FK playlist\_id\\ \FK track\_id\\ position (DEFERRABLE)\\ added\_at, updated\_at}
    \node[umlw] (wm) at (15.2,0) {\textcolor{white}{\rmfamily\bfseries ETL\_WATERMARK}\nodepart{two}\PK id\\ entity\_name (UK)\\ watermark\_column\\ last\_value\\ last\_run\_at\\ rows\_loaded, status};

    % ---- fila inferior ----
    \umlclass{user}{(0,-4.3)}{"USER"}{ink}{\PK id (BIGINT)\\ username, email\\ password, app\_role\\ country, birth\_date\\ gender, ethnicity\\ created\_at, updated\_at}
    \umlclass{ord}{(3.8,-4.3)}{"ORDER"}{ink}{\PK id (BIGINT)\\ \FK user\_id\\ total\_amount, currency\\ payment\_method\\ status (COMPLETED)\\ created\_at, updated\_at}
    \umlclass{oit}{(7.6,-4.3)}{ORDER\_ITEM}{ink}{\PK id (BIGINT)\\ \FK order\_id\\ \FK track\_id\\ price\_paid, currency\\ created\_at, updated\_at}
    \umlclass{pl}{(11.4,-4.3)}{PLAYLIST}{accent}{\PK id (BIGINT)\\ \FK owner\_id\\ name, description\\ tags[], cover\_image\\ is\_public\\ created\_at, updated\_at}
    \node[umlo] (tel) at (15.2,-4.3) {\textcolor{white}{\rmfamily\bfseries TELEMETRY\_EVENT}\nodepart{two}\PK id (offset outbox)\\ event\_id (UUID UK)\\ event\_type, ts\\ session\_id, \FK user\_id\\ device, geo, entrypoint\\ payload (JSONB)\\ created\_at, acked\_at};

    % ---- relaciones principales ----
    \draw[rel] (user.north) -- node[verb,midway,left]{0..1} (artist.south);
    \draw[rel] (artist.east) -- (track.west);
    \draw[rel] (track.east) -- (genre.west);
    \draw[rel] (genre.east) -- (plt.west);
    \draw[rel] (user.east) -- (ord.west);
    \draw[rel] (ord.east) -- (oit.west);
    \draw[rel] (track.south) -- (oit.north);
    \draw[rel] (plt.south) -- (pl.north);
  \end{tikzpicture}
  \end{adjustbox}

  \vspace{0.25em}
  {\tiny\color{muted}
  \textcolor{ink}{\texttt{trig: updated\_at}}: 8 triggers en PostgreSQL actualizan \texttt{updated\_at} en cada \texttt{UPDATE} \quad
  \textcolor{wmgreen}{\rule{2mm}{2mm}}\ \texttt{ETL\_WATERMARK}: control incremental de Glue \quad
  \textcolor{accent}{\rule{2mm}{2mm}}\ \texttt{TELEMETRY\_EVENT}: outbox drenado por HTTP hacia Kafka}
\end{frame}

% ===================== 7 · ARQUITECTURA GENERAL =====================
\section{Arquitectura en AWS}
\begin{frame}{Flujo real desplegado en AWS (\texttt{us-east-1})}
  \centering
  \begin{adjustbox}{max width=\textwidth,max totalheight=0.8\textheight,center}
  \begin{tikzpicture}[font=\scriptsize]
    \foreach \x/\lab in {0/Portal y Estado (VPC),3.8/Ingesta (Stream \& Batch),8.0/Data Lake en Amazon S3,12.4/Analítica y ML}
      \node[font=\scriptsize\bfseries,text=muted] at (\x,3.0) {\lab};

    % 1 · portal
    \node[card,text width=25mm,align=center] (web) at (0,1.6)
      {\textbf{Django tracks4u}\\ {\color{muted}EC2 $\cdot$ S3 Media (.opus)\\ Outbox \texttt{/api/events/}}};
    \node[cylinder,shape border rotate=90,aspect=0.22,draw=ink,fill=white,line width=0.7pt,
          minimum width=18mm,minimum height=11mm,align=center,font=\tiny] (db) at (0,-0.5) {RDS PostgreSQL\\ \texttt{tracks4u}};
    \begin{scope}[on background layer]
      \node[draw=line,dashed,rounded corners=6pt,fit=(web)(db),inner sep=2.5mm] (portal) {};
    \end{scope}

    % 2 · ingesta
    \node[card,text width=26mm,align=center,draw=accent,fill=accent!7] (broker) at (3.8,1.6)
      {\textbf{\color{accent}Kafka en EC2}\\ {\color{muted}\texttt{producer.py} (\texttt{/drain/})\\ Topic \texttt{tracks4u.events}\\ \texttt{consumer.py} (dedup)}};
    \node[card,text width=26mm,align=center,draw=batch,fill=batch!7] (ext) at (3.8,-0.5)
      {\textbf{\color{batch}AWS Glue + Step Fn}\\ {\color{muted}JDBC en VPC $\cdot$ SSL\\ \texttt{etl\_watermark}}};

    % 3 · data lake
    \node[cardsoft,text width=34mm,minimum height=34mm] (lake) at (8.0,0.65) {};
    \node[card,text width=29mm,align=center,inner sep=3pt] (raw) at (8.0,1.5) {\textbf{RAW} {\color{muted}$\cdot$ CSV particionado}};
    \node[card,text width=29mm,align=center,inner sep=3pt] (cur) at (8.0,0.65) {\textbf{CURATED} {\color{muted}$\cdot$ Parquet limpio}};
    \node[card,text width=29mm,align=center,inner sep=3pt,draw=accent,fill=accent!7] (ana) at (8.0,-0.2) {\textbf{\color{accent}ANALYTICS} {\color{muted}$\cdot$ Parquet ML}};
    \node[text=muted,font=\tiny] at (8.0,-0.82) {\texttt{s3://tracks4u-datalake-463790652932-us-east-1}};

    % 4 · consumo
    \node[card,text width=27mm,align=center] (dw) at (12.4,1.6)
      {\textbf{Dashboard Charts}\\ {\color{muted}Ventas por día\\ Ingresos por género\\ Top tracks}};
    \node[card,text width=27mm,align=center,draw=accent,fill=accent!7] (ml) at (12.4,-0.2)
      {\textbf{\color{accent}\texttt{recommend.py}}\\ {\color{muted}Item-Item Coseno\\ \texttt{recommendations}\\ \texttt{user\_profile} $\cdot$ \texttt{metrics}}};

    % flechas
    \draw[flow] (web) -- (db);
    \draw[flowS] (web.east) -- node[above,font=\tiny,text=accent]{drain + ack} (broker.west);
    \draw[flowB] (db.east) -- node[below,font=\tiny,text=batch]{SELECT} (ext.west);
    \draw[flowS] (broker.east) -- (broker.east -| raw.west);
    \draw[flowB] (ext.east) -- (5.9,-0.5) |- ([yshift=-1.2mm]raw.west);
    \draw[flow] (raw) -- (cur);
    \draw[flowS] (raw.east) -- ++(0.35,0) |- (ml.west);
    \draw[flow,{Stealth[length=2mm]}-{Stealth[length=2mm]}] (ana.east) -- (ml.west);
    \draw[flow] (db.south) -- ++(0,-0.55) -| (dw.south);

    \node[cardsoft,text width=132mm,align=center,font=\tiny] at (6.3,-2.35)
      {\textbf{Seguridad e Infraestructura:} IAM Role (\texttt{LabRole}) en EC2 sin credenciales estáticas $\cdot$ RDS en subred privada con \texttt{sslmode=require} $\cdot$ S3 privado con URLs firmadas};
  \end{tikzpicture}
  \end{adjustbox}
\end{frame}

% ===================== 8 · ZONAS DEL DATA LAKE =====================
\section{Arquitectura del data lake}
\begin{frame}{Zonas y particionado en \texttt{tracks4u-datalake-463790652932-us-east-1}}
  \centering
  \begin{adjustbox}{max width=\textwidth,max totalheight=0.8\textheight,center}
  \begin{tikzpicture}[font=\scriptsize]
    \tikzset{zone/.style={card,text width=38mm,align=left,minimum height=30mm,anchor=north},
             chip/.style={draw=line,fill=soft,rounded corners=3pt,inner sep=3pt,font=\scriptsize\ttfamily}}

    \node[zone] (raw) at (0,0) {%
      \textcolor{rawc}{\bfseries RAW}\\ \textbf{Datos crudos (CSV)}\\[3pt]
      \textcolor{muted}{$\bullet$ \texttt{entity=play\_events} (Kafka)}\\
      \textcolor{muted}{$\bullet$ Tablas de negocio (Glue)}\\
      \textcolor{muted}{$\bullet$ Deduplicado en consumer}\\
      \textcolor{muted}{$\bullet$ Commit Kafka post-S3}};
    \node[zone] (cur) at (5.0,0) {%
      \textcolor{batch}{\bfseries CURATED}\\ \textbf{Datos limpios (Parquet)}\\[3pt]
      \textcolor{muted}{$\bullet$ Tipado y validado}\\
      \textcolor{muted}{$\bullet$ Compresión columnar}\\
      \textcolor{muted}{$\bullet$ Cruce catálogo + órdenes}};
    \node[zone,draw=accent,fill=accent!5] (ana) at (10.0,0) {%
      \textcolor{accent}{\bfseries ANALYTICS}\\ \textbf{Salidas ML (Parquet)}\\[3pt]
      \textcolor{muted}{$\bullet$ \texttt{entity=recommendations}}\\
      \textcolor{muted}{$\bullet$ \texttt{entity=user\_profile}}\\
      \textcolor{muted}{$\bullet$ \texttt{entity=model\_metrics}}\\
      \textcolor{muted}{$\bullet$ Top-10 por usuario + \texttt{skip\_rate}}};

    \coordinate (zy) at (0,-1.3);
    \draw[flow] (raw.east |- zy) -- (cur.west |- zy);
    \draw[flow] (cur.east |- zy) -- (ana.west |- zy);

    % jerarquía de particiones
    \node[font=\scriptsize\bfseries,anchor=west] at (-2.1,-3.55) {Particionado estándar en S3: zona / entidad / fecha de ingesta};
    \node[chip,anchor=west] (c1) at (-2.1,-4.25) {s3://tracks4u-datalake-.../};
    \node[chip,right=4mm of c1] (c2) {<zona>/};
    \node[chip,right=4mm of c2,draw=accent,text=accent] (c3) {entity=<entidad>/};
    \node[chip,right=4mm of c3,draw=batch,text=batch] (c4) {ingest\_date=AAAA-MM-DD/};
    \node[chip,right=4mm of c4] (c5) {part-0000.parquet};
    \draw[flow] (c1) -- (c2); \draw[flow] (c2) -- (c3); \draw[flow] (c3) -- (c4); \draw[flow] (c4) -- (c5);

    % ejemplos de rutas reales
    \node[cardsoft,anchor=north west,text width=139mm,align=left,font=\tiny\ttfamily] at (-2.1,-4.95) {%
      raw/\textcolor{accent}{entity=play\_events}/\textcolor{batch}{ingest\_date=2026-11-18}/events\_1763495000000.csv\\
      analytics/\textcolor{accent}{entity=recommendations}/\textcolor{batch}{ingest\_date=2026-11-18}/part-0000.parquet \quad (2\,000 filas: Top-10 $\times$ 200 usuarios)\\
      analytics/\textcolor{accent}{entity=user\_profile}/\textcolor{batch}{ingest\_date=2026-11-18}/part-0000.parquet \qquad (200 filas: conteos + skip\_rate)\\
      analytics/\textcolor{accent}{entity=model\_metrics}/\textcolor{batch}{ingest\_date=2026-11-18}/part-0000.parquet \qquad (Hit-Rate@10 modelo vs popular vs azar)};
  \end{tikzpicture}
  \end{adjustbox}
\end{frame}

% ===================== 9 · MODELO DE RECOMENDACIÓN Y RESULTADOS =====================
\section{Machine Learning en ANALYTICS}
\begin{frame}{\pillM\ \ Recomendador Colaborativo Item-Item y Evaluación}
  \footnotesize
  \begin{columns}[T,onlytextwidth]
    \begin{column}{0.52\textwidth}
      \textbf{\color{accent}Cómo funciona \texttt{recommend.py}}
      \begin{itemize}\setlength{\itemsep}{2pt}
        \item \textbf{Matriz de afinidad ($M$):} agrega eventos de \texttt{raw/entity=play\_events/} ponderados por intención ($+1$ play, $+3$ complete, $+4$ like/playlist, $+5$ purchase, $-2$ skip) con $\text{clip}(M \ge 0)$.
        \item \textbf{Similitud Item-Item ($S$):} similitud coseno entre vectores de canciones ($M^T$), puntaje $R = M \cdot \text{sim}(M^T)$ y enmascarado ($-\infty$) de canciones ya escuchadas.
        \item \textbf{Limpieza híbrida:} extrae \texttt{track\_id} desde \texttt{payload} JSON, unifica IDs reales y sintéticos y descarta anónimos.
        \item \textbf{Protocolo Leave-One-Out:} oculta 1 canción positiva por usuario y evalúa \textbf{Hit-Rate@10}.
      \end{itemize}
    \end{column}
    \begin{column}{0.45\textwidth}
      \textbf{\color{accent}Resultados reales (9\,000 eventos, 200 usuarios, 300 canciones)}
      \vspace{0.3em}

      \renewcommand{\arraystretch}{1.25}
      \begin{tabularx}{\textwidth}{@{}X r@{}}
        \toprule
        \textcolor{muted}{Método evaluado} & \textcolor{muted}{Hit-Rate@10}\\
        \midrule
        \textbf{\color{accent}Modelo Item-Item (Coseno)} & \textbf{\color{accent}0.205 (20.5\%)}\\
        Baseline Más Populares & 0.030 (3.0\%)\\
        Referencia Aleatoria ($10/300$) & 0.033 (3.3\%)\\
        \bottomrule
      \end{tabularx}

      \vspace{0.5em}
      {\scriptsize\color{muted}El modelo Item-Item supera en \textbf{$6.8\times$} al ranking de popularidad global al capturar los clusters de género de cada oyente.}

      \vspace{0.4em}
      \textbf{\color{batch}\scriptsize Entregables en \texttt{s3://.../analytics/}}
      {\tiny\ttfamily\par
      $\bullet$ entity=recommendations (user\_id, rank, track\_id, score)\par
      $\bullet$ entity=user\_profile (conteos por evento + skip\_rate)\par
      $\bullet$ entity=model\_metrics (hit\_rate\_modelo, popular, azar)}
    \end{column}
  \end{columns}
\end{frame}

% ===================== 10 · CIERRE =====================
\begin{frame}[plain]
  \vfill
  \kick{De punta a punta en AWS}
  {\color{ink}\LARGE tracks4u captura cada interacción en Django y RDS,\\ la transporta por Glue y Kafka hacia Amazon S3,\\ y la convierte en recomendaciones con Machine Learning.}\par
  \vspace{0.8em}{\color{accent}\rule{3cm}{2.5pt}}
  \vfill
\end{frame}

\end{document}
